\documentclass[11pt]{article}
\usepackage[utf8]{inputenc}
\usepackage{amsmath,amssymb}
\usepackage[margin=1in]{geometry}
\usepackage{hyperref}
\emergencystretch=3em

\title{Headline part X: the stochastic $1/\log N$ regime}
\author{Claude, with Daniel Murfet}
\date{2026-09-07}

\begin{document}
\maketitle
\sloppy

\begin{abstract}
Paper-facing wrapper for unit~194: at chart level ($d=2$, equal starting exponents), for a
deterministic corner-vanishing observable amplitude, a positive normaliser amplitude and a common
random phase with Taylor data $X_n\Rightarrow Z$,
$\log N_n\cdot Z_{N_n}[\phi](X_n)/Z_{N_n}[1](X_n)\Rightarrow B^\phi_p(Z)/A^1_p(Z)$
(\texttt{headline\_posterior\_log\_decay\_stochastic}), stated on the explicit chart integrals
\texttt{twoDAmp} with the analytic amplitudes of the Taylor data. This is the stochastic form of
\texttt{headline\_posterior\_log\_decay}; when $B^\phi_p(Z)=0$ the limit is $0$ and no faster rate is
claimed. Zero \texttt{sorry}/\texttt{axiom}.
\end{abstract}

\section{Paper-facing meaning}
The empirical posterior expectation of an observable vanishing at the corner of the chart decays like
$1/\log N$ with a \emph{random} coefficient determined by the limiting Taylor data of the empirical
phase; this closes the ``stochastic version of the $1/\log N$ regime'' left open in
\texttt{rem:empirical\_expectation\_lean} since unit~169.

\section{Lean notes}
The wrapper is a one-line \texttt{exact}: \texttt{chartZ} unfolds definitionally to the
\texttt{twoDAmp} of the analytic amplitude of the coefficient pair.

\end{document}
